\section{Continuities and Contrasts with Pre-Christian Writings}
\label{app:pre-christian}

The Church Fathers read Israel's Greek Scriptures alongside Second Temple Jewish writings and the poets and philosophers of Greece. Extra-canonical Jewish works are graded \textbf{Apocryphal}; pagan philosophical systems are received where they touch created spiritual reality and corrected or rejected wherever they confuse creatures with deities, assign evil to nature, or substitute demonic mediation and theurgy for Christ. Where an ancient text is treated through patristic quotation or summary records rather than a standalone complete edition, that transmission basis is stated directly.

{\footnotesize
\setlength{\tabcolsep}{3pt}
\renewcommand{\arraystretch}{1.15}
\begin{longtable}{>{\bfseries\raggedright\arraybackslash}p{0.17\linewidth}>{\raggedright\arraybackslash}p{0.14\linewidth}>{\raggedright\arraybackslash}p{0.22\linewidth}>{\raggedright\arraybackslash}p{0.23\linewidth}>{\raggedright\arraybackslash}p{0.17\linewidth}}
\toprule
\textbf{Source / Corpus} & \textbf{Standing} & \textbf{What the Fathers Received} & \textbf{What the Fathers Corrected or Rejected} & \textbf{Patristic Controls} \\
\midrule
\endfirsthead
\toprule
\textbf{Source / Corpus} & \textbf{Standing} & \textbf{What the Fathers Received} & \textbf{What the Fathers Corrected or Rejected} & \textbf{Patristic Controls} \\
\midrule
\endhead
\bottomrule
\endfoot
Septuagint (LXX): Deut~32:8--9 & Canonical OT (Greek text) & Territorial angels of the nations (\gk{angel\={o}n theou}) under one Providence, paired with Dan~10:13, 20--21. & Rejected any polytheistic division of sovereignty or angelic failure as the source of Gentile idolatry. & Pseudo-Dionysius, \CH{9.2--4}; Jerome, \emph{In Dan.}~10; Basil. \\
\addlinespace[2pt]
Septuagint: Ps~8:6, 95:5, 96:7, 103:4; Job~38:7 & Canonical OT (Greek text) & Angels as spiritual ministers (\gk{pneumata}); Christ worshipped by all angels (Heb~1:6); pagan deities identified as \gk{daimonia} (Ps~95:5); angels praising at creation (Job~38:7). & Guarded against Arian confusion of the Son with angels and against subtle-body materialism. & Athanasius; Basil, \emph{Hex.}~I.5; Augustine, \emph{Civ.\ Dei}~IX.19, XI.9. \\
\addlinespace[2pt]
Septuagint: Is~9:5 [MT/Vulg.~9:6] & Canonical OT (Greek text) & ``Messenger of great counsel'' (\gk{megal\={e}s boul\={e}s angelos}) applied to Christ's mission announcing the Father's counsel. & Corrected Arian attempts to deduce a created angelic nature in the Son; \gk{angelos} signifies mission, not creaturehood. & Justin, \emph{Dial.}~56--61; Athanasius; Pseudo-Dionysius, \CH{4.4}. \\
\addlinespace[2pt]
\emph{1~Enoch} (Watchers; Ch.~20 archangels; prophecy in Jude~1:14--15) & \textbf{Apocryphal} (patristic control; Jude~1:14--15 canonical) & Canonical citation of Enoch's prophecy of judgment in Jude~1:14--15; early 2nd-century use of the Watchers narrative against pagan idolatry (Justin, \emph{2~Apol.}~5). & Decisively rejected carnal union of angels with women in Gen~6:1--4, denied canonical authority to Enochic compilations, and barred unscriptural archangel names from liturgical invocation. & Justin, \emph{2~Apol.}~5 (early reception); Augustine, \emph{Civ.\ Dei}~XV.23 (decisive rejection); Roman Synod~745. \\
\addlinespace[2pt]
\emph{Book of Jubilees} (Ch.~10 Mastema; Day-1 tradition) & \textbf{Apocryphal} (Ch.~10 passage record) & Theological recognition that evil spirits act solely under God's permitting decree (cf.~Job~1--2; Luke~8:31). & Rejected Watchers-offspring etiology and quasi-independent Mastema; Day-1 angelic creation derived strictly from Gen~1:1--5 exegesis rather than apocryphal authority. & Augustine, \emph{Civ.\ Dei}~XI.9, XV.23; Gregory, \emph{Moralia}~XXXII. \\
\addlinespace[2pt]
\emph{Testaments of the Twelve Patriarchs} & \textbf{Apocryphal} (patristic summary) & Early non-canonical witness to Jewish expectation and spiritual combat against spirits of deceit. & Anchored the nine choirs strictly in canonical OT and Pauline texts rather than apocryphal multi-heaven ascents. & Pseudo-Dionysius, \CH{6--9}; Nicene and post-Nicene Fathers. \\
\addlinespace[2pt]
Philo of Alexandria (\emph{De gig.}, \emph{De somn.}, \emph{De conf.\ ling.}, \emph{De opif.}) & Hellenistic Jewish philosophy (patristic summary) & Incorporeal conception of angels as ministers of the Logos and spiritual exegesis of Jacob's ladder and biblical theophanies. & Rejected Philo's identification of angels with disembodied human souls, angelic co-creation of matter or man, and neutral equation of angels with \gk{daimones}. & Irenaeus, \emph{Adv.\ haer.}~II.30.6, IV.7.4, IV.20.1; Augustine, \emph{Civ.\ Dei}~IX.19, XII.24. \\
\addlinespace[2pt]
Hesiod (\emph{Works and Days}) \& Plato (\emph{Symp.}~202d--203a; \emph{Timaeus}) & Greek poetry and philosophy (\emph{Symp.} passage records; Augustine) & Existence of rational immortal spirits created by the supreme God and blessed by cleaving to Him (\emph{Civ.\ Dei}~IX.23 on \emph{Timaeus}). & Rejected golden-age hero \gk{daimones} as mediators and rejected Diotima's intermediary \gk{daimon} in favor of Christ the sole Mediator (1~Tim~2:5). & Augustine, \emph{Civ.\ Dei}~VIII.14, IX.15, IX.19, IX.23. \\
\addlinespace[2pt]
Plutarch \& Apuleius (\emph{De deo Socratis}) & Middle Platonism (Apuleius quoted verbatim in Augustine) & Recognized the distinction between immortal and mortal rational creatures while exposing pagan religious longing for mediation. & Refuted Apuleius's fivefold definition (passible in soul, aerial in body, eternal in time): a passible, angry spirit is worse than a righteous mortal man and unfit to mediate. & Augustine, \emph{Civ.\ Dei}~VIII.14--22, IX.15. \\
\addlinespace[2pt]
Plotinus, Porphyry (\emph{Letter to Anebo}), \& Iamblichus (\emph{De mysteriis}) & Neoplatonism (Plotinus and Porphyry quoted and answered in Augustine) & Plotinus's insight on mortal mercy vs.\ eternal demonic misery (\emph{Civ.\ Dei}~IX.10); Porphyry's maxim that holy angels seek imitation rather than sacrifice (X.26); triadic participation purified by Dionysius. & Rejected necessary emanationism in favor of creation \emph{ex nihilo}; exposed Iamblichean and Porphyrian theurgy as veiled demonolatry (X.9--11). & Augustine, \emph{Civ.\ Dei}~IX.10, X.9--11, X.26; Pseudo-Dionysius, \CH{1--4}, \emph{DN}~IV.23. \\
\end{longtable}
}
